\scanpage{53}
\mspage{52}

\textcircled{3} Cas des ateliers spacieux ou polyédraux, axiomatique en termes de $(\unsure{M|\mathcal M},\leq,\ll,\parallel_{\unsure{\mathcal L|\mathscr L}})$, où $\parallel_{\mathcal L}$ est une relation ``de disjonction'' sur $\mathcal L={}\text{ens. des él. min. de }(M,\ll)$. (Cf p. 34)

\margin{[Prat dis 1 à Prat dis 5 comme Prat $\unsure{\Lambda|\wedge}$ 1 à Prat $\Lambda$ 5, avec $\Lambda=\mathcal L$]}

% slanted
\margin{dis pour ``disjonction'' ``Prat'' pour ``préatelier''}

\[
\text{\uline{Prat \deleted{\unsure{dis}}\inserted{$\mathcal L$} 1$'$}}\quad X\leq Y\Rightarrow X\ll Y
\]

\[
\text{\uline{Prat \deleted{dis}\inserted{$\mathcal L$} 2$'$}}\quad \text{\deleted{\unsure{Comme}} }\forall X\in M,\quad \mathrm{omb}(X)^{\circ}\ne\varnothing\quad \text{i.e. }\exists x\in\mathcal L,\ x\overset{\circ}{\ll}X
\]

\[
\text{\uline{Prat \deleted{dis}\inserted{$\mathcal L$} 3$'$}}\quad \forall X\in M,\ x\in\mathcal L,\ x\ll X,\quad \exists X'\in M,\quad x\overset{\circ}{\ll}X'\leq X\unsure{,|.}
\]

\uline{Prat \deleted{dis}\inserted{$\mathcal L$} 4$'$} Comme Prat $\deleted{\unsure{A}}\Lambda$ 4 : \unsure{i.e.} Soit

\diagram{
\[
\begin{aligned}
&Y\ll\vcenter{\hbox{\begin{tikzcd}[ampersand replacement=\&,row sep=0pt,cells={nodes={inner sep=0pt}}]
X\\
\unsure{\rotatebox{90}{$\leq$}|\vee\! /}\\
X'
\end{tikzcd}}}\\
&\text{et }x\in\mathcal L,\quad x\ll Y,X',\quad \text{alors }\exists Y'\ \text{avec}\quad Y'\leq Y,\ Y'\ll X',\ x\ll Y'\unsure{,|.}
\end{aligned}
\]
}

\[
\text{i.e.}\qquad \mathrm{omb}(Y)\cap\mathrm{omb}(X')=
\underset{\substack{\parallel\,(\mathrm{déf})\\ \displaystyle\bigcup_{Z\in Y_{X'}}\mathrm{omb}(Z)}}{\mathrm{omb}(Y_{X'})}
\quad \text{où }Y_{X'}=\widetilde Y\cap\mathrm{omb}(X'),
\]

De plus, si $Y\leq X$, \unsure{OPS|ops} $Y'\leq X$, i.e

\[
\mathrm{omb}(Y)\cap\unsure{\mathrm{om}|\mathrm{omb}}(X')=\mathrm{omb}(L)\quad \text{où}\quad L=\widetilde Y\cap\widetilde{X'}.
\]

\uline{Prat \deleted{dis}\inserted{$\mathcal L$} 5$'$} $\parallel_{\mathcal L}$ est \deleted{réflexive}\inserted{symétrique} et antiréflexive (i.e. $x\deleted{=}\parallel y$ implique $y\parallel x$ et $x\ne y$).

\[
\begin{aligned}
&\text{\uline{Prat \deleted{dis}\inserted{$\mathcal L$} 6$'$}}\quad \forall X,Y\in M\ \text{avec }Y\leq X,\quad \deleted{\unsure{xy|x,y}}\ x,y\in\mathcal L,\ y\ll Y,\\
&x\ll X\ \text{et }\deleted{y\in\mathrm{omb}(\unsure{X|Y}}\ x\not\ll Y,\quad (\text{i.e. } y\in\mathrm{omb}(Y),x\in\mathrm{omb}(X)\setminus\mathrm{omb}(Y))\\
&\text{on a }\ x\parallel y.
\end{aligned}
\]

\margin{[Seuls les deux derniers axiomes font intervenir $\parallel$]}

Ce sont les \unsure{$\overline{\mathrm{ms}}$|$\Longleftrightarrow$} axiomes que Prat $\Lambda$, à l'exception du dernier Prat dis 6 (et à l'exception du fait que Prat dis 5 prend une forme plus simple) -- et qui \unsure{remplacent|remplace} \deleted{\unsure{les}} \unsure{douzaines|dizaines|dizaine|derniers} d'axiomes, \unsure{correspondant|correspond.t|en exprimant} \unsure{$\mathrel{>\mkern-13mu<}$|et} en termes d'

% slanted
\margin{Donc le \uline{\unsure{$\overline{\mathrm{m}}$|$\overline{\mathrm{me}}$}} préatelier, permet de construire \unsure{à la fois|à la|à} un atelier polyédral et un atelier spacieux (le plus \unsure{souvent|souvt|svt} différ\unsure{ents|ts}\ldots)}

On définira :

\[
X\mathrel{>\mkern-13mu<}Y\Longleftrightarrow
\left\{\begin{aligned}
&\forall\deleted{\unsure{a}}x\in\mathrm{omb}(X)\setminus\mathrm{omb}(\deleted{\unsure{X}}L),y\in\mathrm{omb}(Y)\setminus\mathrm{omb}(L)\\
&\qquad(\text{où }L=\widetilde X\cap\widetilde Y)\ \text{on a }\ x\parallel y\\
&[+\ L=\varnothing\ \text{ou } L\ \text{de la forme }\widetilde Z,\ \text{dans le}\\
&\hspace{16em}\text{cas ``polyédral''}\ ]
\end{aligned}\right.
\]

